% ---
% title: Venn diagrams and inclusion–exclusion
% tags: venn-diagrams, inclusion-exclusion
% ---

\begin{problem}
	A town of 100\,000 inhabitants has three newspapers I, II, III. The
	percentages of readers: I: 10, II: 30, III: 5, I and II: 8, I and III: 2,
	II and III: 4, all three: 1. I and III are morning papers, II is an
	evening paper. How many people read (a) exactly one paper? (b) at least two
	papers? (c) at least one morning paper and the evening paper? (d) no paper?
	(e) exactly one morning paper and the evening paper?
\end{problem}

\begin{slideproblem}
	A town of 100\,000 has newspapers I, II, III. Readers (\%):
	I: 10, II: 30, III: 5, I\&II: 8, I\&III: 2, II\&III: 4, all: 1.
	How many read (a) exactly one paper? (b) at least two?
	(c) a morning paper (I or III) and the evening paper II?
	(d) no paper? (e) exactly one morning paper and the evening paper?
\end{slideproblem}

\begin{slides}
	\begin{columns}[T]
		\begin{column}{.6\textwidth}
			\small
			\exstatement
			\uncover<2->{\sol Fill the Venn diagram \textbf{from the inside out}:
					\begin{itemize}
						\item center (all three): $1$;
						\item only I\&II: $8-1=7$; only I\&III: $2-1=1$; only II\&III:
						      $4-1=3$;
						\item only I: $10-7-1-1=1$; only II: $30-7-3-1=19$; only III:
						      $5-1-3-1=0$.
					\end{itemize}}
		\end{column}
		\begin{column}{.38\textwidth}
			\centering
			\begin{tikzpicture}[scale=.6]
				\draw[navy,thick] (0,0) circle (1.5);
				\draw[teal,thick] (1.8,0) circle (1.5);
				\draw[gold,thick] (.9,-1.5) circle (1.5);
				\uncover<2->{
						\node at (-.7,.3) {\scriptsize 1};
						\node at (2.5,.3) {\scriptsize 19};
						\node at (.9,-2.3) {\scriptsize 0};
						\node at (.9,.6) {\scriptsize 7};
						\node at (.1,-.95) {\scriptsize 1};
						\node at (1.7,-.95) {\scriptsize 3};
						\node at (.9,-.45) {\scriptsize 1};}
				\node[navy] at (-1.4,1.4) {I};
				\node[teal] at (3.2,1.4) {II};
				\node[gold] at (.9,-3.3) {III};
			\end{tikzpicture}\\
			\footnotesize (in \%)
		\end{column}
	\end{columns}
\end{slides}

\begin{slides}
	\small
	\sol (percentages of 100\,000 people; read off the diagram)
	\begin{enumerate}
		\item[(a)] Exactly one: only I $+$ only II $+$ only III
		      $=1+19+0=20\%$: \textbf{20\,000}.
		\item[(b)] At least two: the four inner regions $7+1+3+1=12\%$:
		      \textbf{12\,000}.
		\item[(c)] $(I\cup III)\cap II = (I\cap II)\cup(III\cap II)$; the two
		      overlap in ``all three'': $8+4-1=11\%$: \textbf{11\,000}.
		\item[(d)] Inclusion--exclusion:
		      \[ |I\cup II\cup III| = 10+30+5-8-2-4+1 = 32\%, \]
		      so $100-32=68\%$ read none: \textbf{68\,000}.
		\item[(e)] II together with exactly one of I, III: only I\&II $+$ only
		      II\&III $=7+3=10\%$: \textbf{10\,000}.
	\end{enumerate}
	\begin{block}{Inclusion--exclusion for three sets}
		$|E\cup F\cup G| = |E|+|F|+|G|-|EF|-|EG|-|FG|+|EFG|$
	\end{block}
\end{slides}

\begin{solution}
	\idea The given numbers overlap (``I and II'' includes those who also read
	III). Fill the Venn diagram \emph{from the inside out}; then every region
	holds the people of exactly that combination, and we can simply add regions.

	\begin{minipage}{.58\linewidth}
		\step{Filling the diagram (in \%).}
		\begin{itemize}[nosep]
			\item center (all three): $1$;
			\item only I\&II: $8-1=7$; only I\&III: $2-1=1$; only II\&III: $4-1=3$;
			\item only I: $10-7-1-1=1$;\\ only II: $30-7-3-1=19$;\\ only III:
			      $5-1-3-1=0$.
		\end{itemize}
	\end{minipage}\hfill
	\begin{minipage}{.38\linewidth}
		\centering
		\begin{tikzpicture}[scale=.65]
			\draw[navy,thick] (0,0) circle (1.5);
			\draw[teal,thick] (1.8,0) circle (1.5);
			\draw[gold,thick] (.9,-1.5) circle (1.5);
			\node at (-.7,.3) {\small 1};
			\node at (2.5,.3) {\small 19};
			\node at (.9,-2.3) {\small 0};
			\node at (.9,.6) {\small 7};
			\node at (.1,-.95) {\small 1};
			\node at (1.7,-.95) {\small 3};
			\node at (.9,-.45) {\small 1};
			\node[navy] at (-1.4,1.4) {I};
			\node[teal] at (3.2,1.4) {II};
			\node[gold] at (.9,-3.3) {III};
		\end{tikzpicture}
	\end{minipage}

	\step{(a)} Exactly one: $1+19+0=20\%$, i.e.\ $20\,000$ people.
	\step{(b)} At least two: the four inner regions $7+1+3+1=12\%$: $12\,000$.
	\step{(c)}
	\[ (\mathrm{I}\cup\mathrm{III})\cap\mathrm{II}=
		(\mathrm{I}\cap\mathrm{II})\cup(\mathrm{III}\cap\mathrm{II}). \]
	Inclusion--exclusion:
	$8+4-1=11\%$ (the readers of all three were counted twice). Or from the
	diagram: $7+3+1=11\%$: $11\,000$.
	\step{(d)}
	\[ |\mathrm{I}\cup\mathrm{II}\cup\mathrm{III}|=10+30+5-8-2-4+1=32\% \]
	(check: the seven regions sum to $1+19+0+7+1+3+1=32$), so $100-32=68\%$ read
	none: $68\,000$.
	\step{(e)} II together with exactly one of I, III (not both): only I\&II
	plus only II\&III: $7+3=10\%$: $10\,000$.
	\begin{answer}
		(a) $20\,000$ \quad (b) $12\,000$ \quad (c) $11\,000$ \quad
		(d) $68\,000$ \quad (e) $10\,000$
	\end{answer}
	\begin{remark}
		\textbf{Common mistake:} computing ``only I'' as $10-8-2=0$. This
		subtracts the center (readers of all three) twice. Always start from the
		center.
	\end{remark}
\end{solution}
